% TODO don't use a pmp function to have a closer translation

% FIXME what is needed as input (µSail, contracts, UC)
\chapter{Katamaran}
Katamaran is a framework for verifying \textmu{}Sail programs using \glspl{uc}.  
A \gls{uc} specifies the expected behaviour of a program in terms of preconditions and postconditions, and Katamaran proves that the given \textmu{}Sail implementation satisfies these contracts.
This makes it possible to establish formal correctness and security properties for ISA models, rather than relying only on testing or simulation.~\cite{universal_contract}

Katamaran is built around \textmu{}Sail, a subset of Sail to enable easier verification.
In this setting, the ISA description serves as the executable model, while the \gls{uc} captures the security-relevant properties that the model must satisfy.
This is particularly useful for architectural mechanisms such as memory protection, where one wants to prove that the model enforces the intended isolation guarantees.

\section{The RISC-V PMP Case Study}
Katamaran includes a case study for the RISC-V \gls{pmp} extension.
The case study formalizes the \gls{pmp} behaviour in \textmu{}Sail and proves a \gls{uc} expressing the memory-isolation and access-control properties expected from \gls{pmp}.

The \textmu{}Sail model is derived from the official RISC-V Sail implementation~\cite{risc-v_sail} as you can see in the \texttt{pmpCheckRWX}-function in~\cref{lst:sail_code, lst:microsail_code}. %TODO insert these listings and look for a good example function ¿ pmpCheckRWX ?
As a result, this model stays close to the real-world specification.

The existing model is however intentionally limited to serve as a proof of concept.
It only supports two \gls{pmp} entries, rather than 16 or 64 allowed in the real RISC-V specification.
It also only supports the \gls{tor} addressing mode.

The \textmu{}Sail model, just like the official RISC-V Sail model implements \gls{pmp} using a \texttt{pmpCheck} function.
This function is complex enough that it itself calls multiple functions.

To then implement the \gls{uc} every function needs a definition of what it is supposed to do, which is done by defining contracts for every function as seen in~\cref{lst:microsail_contract}.
These contracts have a precondition which needs to be true for every invocation of the function and a postcondition which also needs to be true after the invocation of the corresponding function.
In the given example the precondition is \top which means this condition is always fulfilled, but the postcondition is a bit more complex.
The postcondition starts of with converting the `L', `A', `X', `W', `R' which are part of the pmpcfg-register into a usable form for this function.
Then it is checked whether the access succeded according to the \textmu{}Sail code, if it did we check whether this is supposed to happen, otherwise the function is regarded as correct.

When all the function's contracts are proven true, the prover can use these contracts to prove the \gls{uc}.

% TODO simplify to more closely represent µSail code
\begin{lstlisting}[caption=Sail code example, label=lst:sail_code, basicstyle=\tiny]
private function pmpCheckRWX(ent : Pmpcfg_ent, access : MemoryAccessType(mem_payload)) -> bool =
  match access {
    Load(Data)                    => ent[R] == 0b1,
    Load(PageTableEntry)          => ent[R] == 0b1,
    LoadReserved(Data)            => ent[R] == 0b1,
    Store(Data)                   => ent[W] == 0b1,
    Store(PageTableEntry)         => ent[W] == 0b1,
    StoreConditional(Data)        => ent[W] == 0b1,
    Atomic(_, Data, Data)         => ent[R] == 0b1 & ent[W] == 0b1,
    InstructionFetch()            => ent[X] == 0b1,

    // "The PMP checks are extended to require read-write permission
    // for memory accessed by shadow stack instructions."
    Load(ShadowStack)             => ent[R] == 0b1 & ent[W] == 0b1,
    Store(ShadowStack)            => ent[R] == 0b1 & ent[W] == 0b1,

    Atomic(_, ShadowStack, ShadowStack) => ent[R] == 0b1 & ent[W] == 0b1,

    LoadReserved(p)               => internal_error(__FILE__, __LINE__, "Invalid payload (" ^ mem_payload_name(p) ^ ") for LoadReserved."),
    StoreConditional(p)           => internal_error(__FILE__, __LINE__, "Invalid payload (" ^ mem_payload_name(p) ^ ") for StoreConditional."),
    Atomic(_, rp, wp)             => internal_error(__FILE__, __LINE__, "Invalid payloads (" ^ mem_payload_name(rp) ^ ", " ^ mem_payload_str(wp) ^ ") for Atomic."),

    // TODO: "If neither a load instruction nor store instruction is
    // permitted to access the physical addresses, but an instruction
    // fetch is permitted to access the physical addresses, whether a
    // cache-block management instruction is permitted to access the
    // cache block is UNSPECIFIED."
    CacheAccess(CB_manage(_))   => ent[R] == 0b1 | ent[W] == 0b1,

    CacheAccess(CB_zero())      => ent[W] == 0b1,
    CacheAccess(CB_prefetch(p)) => match p {
      PREFETCH_I => ent[X] == 0b1,
      PREFETCH_R => ent[R] == 0b1,
      PREFETCH_W => ent[W] == 0b1,
    },
  }
\end{lstlisting}

\begin{lstlisting}[caption=\textmu{}Sail code example, label=lst:microsail_code]
Definition fun_pmpCheckRWX : Stm [ent ∷ ty_pmpcfg_ent; acc ∷ ty_access_type] ty.bool :=
match: ent in rpmpcfg_ent with
  [L; A; X; W; R] =>
    match: acc in union access_type with
    |> KRead pat_unit      => if: R then stm_val ty.bool true else stm_val ty.bool false
    |> KWrite pat_unit     => if: W then stm_val ty.bool true else stm_val ty.bool false
    |> KReadWrite pat_unit => if: R && W then stm_val ty.bool true else stm_val ty.bool false
    |> KExecute pat_unit   => if: X then stm_val ty.bool true else stm_val ty.bool false
    end
end.
\end{lstlisting}

%TODO simplify contract
\begin{lstlisting}[caption=The contract for the \textmu{}Sail code, label=lst:microsail_contract]
Definition sep_contract_pmpCheckRWX : SepContractFun pmpCheckRWX :=
  let Σ : LCtx := [acc :: ty_access_type; L :: ty.bool; A :: ty_pmpaddrmatchtype; X :: ty.bool; W :: ty.bool; R :: ty.bool] in
  let entry : Term Σ _ := term_record rpmpcfg_ent [term_var L; term_var A; term_var X; term_var W; term_var R] in
  {| sep_contract_logic_variables := Σ;
     sep_contract_localstore      := [nenv entry; term_var acc];
     sep_contract_precondition    := ⊤;
     sep_contract_result          := "result_pmpCheckRWX";
     sep_contract_postcondition   :=
       let entry := term_record rpmpcfg_ent [term_var L; term_var A; term_var X; term_var W; term_var R] in
       if: term_var "result_pmpCheckRWX"
       then asn_pmp_check_rwx entry (term_var acc)
       else ⊤;
  |}.



Notation asn_pmp_check_rwx cfg acc := (asn.formula (formula_user pmp_check_rwx [cfg;acc])).
\end{lstlisting}